% TOP_PROBLEM: {"id": 30006861, "title": "Irrationality of zeta(5)", "category_id": 1, "category": {"id": 1, "name": "number_theory", "display_name": "Number Theory", "description": "Properties of integers, prime numbers, Diophantine equations.", "slug": "number-theory", "order_index": 1, "created_at": "2026-07-31T15:26:25.670Z"}, "status": "open"}
% FIRST_POSTED: 2026-09-27
% !TeX program = lualatex
% ATTEMPT_STATUS: unresolved
\documentclass[11pt]{article}
\usepackage[margin=25mm]{geometry}
\usepackage{fontspec}
\setmainfont{Latin Modern Roman}
\usepackage{amsmath,amssymb,mathtools}
\usepackage{unicode-math}
\setmathfont{Latin Modern Math}
\usepackage{xurl,hyperref}
\hypersetup{colorlinks=true,urlcolor=blue}
\setcounter{secnumdepth}{0}
\setlength{\parindent}{0pt}
\setlength{\parskip}{5pt}
\setlength{\emergencystretch}{3em}
\begin{document}
\section*{\texorpdfstring{Irrationality of zeta(5)}{Irrationality of zeta(5)}}\label{30006861}
\subsection{Definitions and mathematical statement}
The absolutely convergent series
\[
 \zeta(5)=\sum_{n=1}^{\infty}\frac1{n^5}
\]
defines a positive real number. The problem is to prove $\zeta(5)\notin{\mathbb{Q}}$, equivalently
\[
 \forall a\in{\mathbb{Z}}\ \forall b\in{\mathbb{Z}}_{\ge1},\qquad b\zeta(5)\ne a.
\]
The target singles out this one value, rather than proving that some member of a collection of odd zeta values is irrational. An approximation with many certified decimal digits does not exclude rationality, since the denominator of a hypothetical rational value is unrestricted. Transcendence, meaning that no nonzero polynomial in ${\mathbb{Q}}[X]$ vanishes at this value, is a stronger assertion and is not requested here.
\par\smallskip\textit{Formulation and concept references:} \href{https://www2.mathematik.hu-berlin.de/~kreimer/wp-content/uploads/IrratModuliMotivesv8.pdf}{[S1]}.

\subsection{Short English statement}
Can the sum of the reciprocals of all positive fifth powers be proved not to be a ratio of integers?
\subsection{Research attempt: integral linear forms, Hankel positivity, and normalization}
\paragraph{1. A current-status qualification that cannot be omitted.}
The target is irrationality of the real number $\zeta(5)$ itself. Brown's source [S1] supplies historical context for integral methods; its supplied URL failed in this audit, and the author's accessible copy [E3] was consulted instead. A primary-source audit on 27 September 2026 also found Fauzan's 17 September 2026 preprint [E1], which claims this exact irrationality result using normalized Hankel determinants. Romik's repository [E2] reports a Lean formalization with a prime-number-theorem axiom as its one nonlogical input. Its README, theorem interface, and axiom declaration were inspected. The formal statement uses the classical real series, not a $p$-adic zeta value. No clean Lean build, dependency audit, or complete verification of the preprint's normalization estimates was performed here. Thus this document records a recent exact-target claim without converting that report into an independently certified proof. Its own mathematical contribution is the explicit analysis below.

\paragraph{2. The precise criterion an approximation must satisfy.}
If there are integers $a_n,b_n$ with
\[
 0<|a_n\xi-b_n|\longrightarrow0,
\]
then $\xi$ is irrational. Indeed, if $\xi=A/B$ with $B>0$, every nonzero $a_n\xi-b_n$ has absolute value at least $1/B$. The coefficients $a_n$ need not be bounded, but the integer linear form must be nonzero and tend to zero after all denominators have been cleared.

A useful polynomial version is equally elementary. Suppose $Q_n\in\mathbb Z[T]$, $\deg Q_n\leq d_n$, and $0<|Q_n(\xi)|$. If for every fixed $B\geq1$ one has $B^{d_n}|Q_n(\xi)|\to0$, then $\xi$ is irrational. At a rational $A/B$, $B^{d_n}Q_n(A/B)$ is a nonzero integer. For example, $d_n=O(n)$ and $|Q_n(\xi)|\leq\exp(-c n^2)$ for a fixed $c>0$ suffice. Degree growth matters: the same bound need not suffice if $d_n$ grows much faster than $n^2$. This is the exact arithmetic conversion needed by determinant methods; numerical smallness before normalization is not enough.

\paragraph{3. A positive integral that isolates the value but not irrationality.}
For an integer $m\geq0$, termwise integration of the nonnegative geometric series gives
\[
 m_m:=\frac1{24}\int_0^1\frac{x^m(-\log x)^4}{1-x}\,dx
 =\sum_{k=m+1}^{\infty}\frac1{k^5}
 =\zeta(5)-H_m^{(5)},
 \qquad H_m^{(5)}=\sum_{k=1}^{m}\frac1{k^5}.
\]
Here $H_0^{(5)}=0$. The elementary identity $\int_0^1x^{k-1}(-\log x)^4dx=24/k^5$ follows by substituting $x=e^{-t}$ and integrating the Gamma integral four times by parts. Monotone convergence justifies the interchange. Integral comparison gives
\[
 \frac1{4(m+1)^4}\leq m_m\leq\frac1{4m^4}\qquad(m\geq1).
\]
If $d_m=\operatorname{lcm}(1,\ldots,m)$, then $d_m^5H_m^{(5)}$ is an integer. But the small tail is only polynomial in $m$ while this natural clearing factor grows exponentially. Therefore this particular integral representation does not meet the integer-linear-form criterion.

Multiplying by an apparently helpful vanishing factor can make matters worse. For $n\geq1$,
\[
 I_n=\frac1{24}\int_0^1x^n(1-x)^{n-1}(-\log x)^4dx
 =\sum_{j=0}^{n-1}\frac{(-1)^j{n-1\choose j}}{(n+j+1)^5}\in\mathbb Q.
\]
This integral is positive, but the factor $(1-x)^n$ in the original moment integrand cancels its pole at one. All zeta terms have disappeared. A rapidly decaying positive rational number cannot prove irrationality of a number that no longer occurs in its expression.

\paragraph{4. A positive Hankel determinant with exact arithmetic structure.}
For $n\geq1$ define
\[
 D_n=\det(m_{i+j})_{0\leq i,j<n},\qquad
 P_n(T)=\det\bigl(T-H_{i+j}^{(5)}\bigr)_{0\leq i,j<n}.
\]
Then $D_n=P_n(\zeta(5))$. Let
\[
 d\mu(x)=\frac{(-\log x)^4}{24(1-x)}\,dx\quad(0<x<1).
\]
This is a finite positive measure with total mass $\zeta(5)$ and positive density throughout the open interval. The Hankel matrix is the Gram matrix of $1,x,\ldots,x^{n-1}$ in $L^2(\mu)$. Every nonzero polynomial has positive squared norm, since it cannot vanish on an interval. The matrix is therefore positive definite, and
\[
 D_n>0.
\]
Equivalently, expanding determinants and applying Fubini gives
\[
 D_n=\frac1{n!}\int_{(0,1)^n}\prod_{a<b}(x_a-x_b)^2\prod_{a=1}^nd\mu(x_a)>0.
\]
The factor $1/n!$ follows because all permutations contribute equally to the squared Vandermonde integral. This representation proves nonvanishing exactly, rather than by floating-point evaluation of an ill-conditioned matrix.

There is a further simplification: $P_n$ has degree at most one, not degree $n$. In the determinant's multilinear expansion by columns, a term choosing $T\mathbf1$ in two columns vanishes because those columns coincide. Only terms with zero or one such column survive. For instance,
\[
 P_1(T)=T,\qquad P_2(T)=\frac{31}{32}T-1.
\]
Positivity for $n=2$ recovers the rational lower bound $\zeta(5)>32/31$. It does not establish irrationality. Positivity also shows that $P_n$ is not the zero polynomial, although its linear coefficient could in principle vanish without contradicting this particular observation.

\paragraph{5. A proved quadratic-exponential upper bound.}
Gram determinants are unchanged by replacing the monomial basis by any monic polynomial basis: the change-of-basis matrix is triangular with diagonal entries one. Let $q_0=1$ and, for $j\geq1$, let
\[
 q_j(x)=2^{1-2j}T_j(2x-1),
\]
where $T_j$ is the Chebyshev polynomial. Since $T_j$ has leading coefficient $2^{j-1}$, $q_j$ is monic, and $|q_j(x)|\leq2^{1-2j}$ for $0\leq x\leq1$. Hadamard's inequality for the Gram matrix yields
\[
 0<D_n\leq\prod_{j=0}^{n-1}\int q_j(x)^2d\mu(x)
 \leq\zeta(5)^n\,2^{-2(n-1)^2}.
\]
The exponent is obtained from $\sum_{j=1}^{n-1}(2-4j)=-2(n-1)^2$. Thus this very simple determinant already has genuine quadratic-exponential decay and exact positivity. The missing ingredient is arithmetic normalization, not real smallness.

To make that point quantitative, set $L_n=\operatorname{lcm}(1,\ldots,2n-2)$ for $n\geq2$. Every $H_{i+j}^{(5)}$ has denominator dividing $L_n^5$. Expanding the determinant gives the safe integral polynomial
\[
 Q_n(T)=L_n^{5n}P_n(T)\in\mathbb Z[T].
\]
The previous estimate alone gives
\[
 0<Q_n(\zeta(5))
 \leq L_n^{5n}\zeta(5)^n2^{-2(n-1)^2}.
\]
Even using the prime-number-theorem asymptotic $\log\operatorname{lcm}(1,\ldots,m)=m+o(m)$, this upper bound has logarithm $(10-2\log2)n^2+o(n^2)$, which is positive in its leading term. It does not tend to zero. This does not prove that the actual normalized values grow; it proves that the established upper bound is inadequate for the criterion. Improved common-divisor extraction or a different rational functional is a substantive missing step, not a harmless rescaling.

\paragraph{6. The normalization issue in a recent claim.}
A rational normalization $c_nP_n$ is useful only if two statements hold simultaneously: all coefficients of $c_nP_n$ are integral, and $c_nP_n(\zeta(5))$ remains sufficiently small. Positivity of the moment representation establishes neither integrality nor a bound for $|c_n|$. Conversely, denominator control alone does not guarantee that the normalized real value decays. These are logically separate estimates.

The current claim [E1] concerns a more elaborate normalized determinant construction; the elementary $P_n$ above is not asserted to reproduce it. The repository [E2] expressly reports dependence on a prime-number-theorem axiom. Since that theorem is established mathematics, this dependency is not by itself an objection to a formal proof. However, reading a README and an axiom declaration is not the same as compiling the project and checking the entire dependency graph and exact target. This attempt has done the former and not the latter. It therefore makes neither a claim that the formal proof is defective nor a claim to have independently certified it.

\paragraph{7. Exact diagnostics and outcome.}
The accompanying rational-arithmetic program constructs $P_n$ for small $n$ by evaluating determinants and recovering the affine polynomial, checks its degree by further exact evaluations, verifies $P_2$, and checks the stated clearing factor. It also verifies the Chebyshev leading coefficients and exponent sum in a finite range. Such checks cannot prove irrationality: they concern finitely many rational polynomials and not the full asymptotic arithmetic normalization.

\textbf{UNRESOLVED as an independent proof.} This attempt establishes exact positive moment forms, a positive Hankel determinant, its affine dependence on $\zeta(5)$, and a quadratic-exponential upper bound. Its own first gap is a normalization with integral coefficients whose cost is smaller than the available decay. Separately, a September 2026 primary-source claim and reported formalization address the exact target; their existence is recorded with explicit verification limits rather than an obsolete unconditional claim that the target is still open.

\subsection{Sources}
\begin{itemize}
\item[S1]Francis Brown, "Irrationality proofs for zeta values, moduli spaces and dinner parties," author-hosted PDF (v8), www2.mathematik.hu-berlin.de.
\url{https://www2.mathematik.hu-berlin.de/~kreimer/wp-content/uploads/IrratModuliMotivesv8.pdf}.
\textit{Formulation source.}
\item[E1]A. Fauzan, $\zeta(5)$ is irrational, Zenodo preprint record, version 1, 17 September 2026.
\url{https://zenodo.org/records/22826419}.
\textit{Primary exact-target claim inspected 27 September 2026. The record's abstract was accessible; its PDF retrieval failed in this audit, and the full proof was not independently certified.}
\item[E2]D. Romik, zeta5-irrationality, Lean formalization repository, README, Zeta5/Interface.lean and Zeta5/Axioms.lean, inspected 27 September 2026.
\url{https://github.com/danromik/zeta5-irrationality}.
\textit{Reports a formal proof using the prime number theorem as one external axiom. No local clean build or complete dependency audit was performed for this attempt.}
\item[E3]F. Brown, Irrationality proofs for zeta values, moduli spaces and dinner parties, author-hosted version 8, Sections 1.6--1.7.
\url{https://www.ihes.fr/~brown/IrratModuliMotivesv8.pdf}.
\textit{Accessible primary copy checked 27 September 2026 after the supplied S1 URL failed; historical integral-method context, not a current resolution claim.}
\end{itemize}
\end{document}
